\begin{ex}
Số liền sau của 999 999 là số nào?
\shortans{1000000}
\loigiai{
1000000.
}
\end{ex}

\begin{ex}
Tính: $450000 + 550000 = \dots$
\shortans{1000000}
\loigiai{
1000000.
}
\end{ex}

\begin{ex}
Tính: $1000000 - 320000 = \dots$
\shortans{680000}
\loigiai{
680000.
}
\end{ex}

\begin{ex}
Tính: $25 \times 40 = \dots$
\shortans{1000}
\loigiai{
1000.
}
\end{ex}

\begin{ex}
Tính: $1200 : 30 = \dots$
\shortans{40}
\loigiai{
40.
}
\end{ex}

\begin{ex}
Rút gọn phân số $\frac{24}{36}$ về tối giản.
\shortans{2/3}
\loigiai{
$\frac{2}{3}$.
}
\end{ex}

\begin{ex}
Tính: $\frac{1}{3} + \frac{1}{2} = \dots$
\shortans{5/6}
\loigiai{
$\frac{5}{6}$.
}
\end{ex}

\begin{ex}
Tính: $\frac{3}{4} - \frac{1}{2} = \dots$
\shortans{1/4}
\loigiai{
$\frac{1}{4}$.
}
\end{ex}

\begin{ex}
Tính: $\frac{2}{5} \times \frac{5}{6} = \dots$ (tối giản).
\shortans{1/3}
\loigiai{
$\frac{1}{3}$.
}
\end{ex}

\begin{ex}
Tính: $\frac{3}{4} : \frac{3}{8} = \dots$
\shortans{2}
\loigiai{
$\frac{3}{4} \times \frac{8}{3} = 2$.
}
\end{ex}

\begin{ex}
Tìm $\frac{2}{3}$ của 60.
\shortans{40}
\loigiai{
$60 \times \frac{2}{3} = 40$.
}
\end{ex}

\begin{ex}
Tìm hai số có tổng là 80 và hiệu là 20. Số lớn là bao nhiêu?
\shortans{50}
\loigiai{
$(80 + 20) : 2 = 50$.
}
\end{ex}

\begin{ex}
Số bé trong bài toán trên là bao nhiêu?
\shortans{30}
\loigiai{
$(80 - 20) : 2 = 30$.
}
\end{ex}

\begin{ex}
Tính trung bình cộng của 3 số: 20, 30, 40.
\shortans{30}
\loigiai{
30.
}
\end{ex}

\begin{ex}
Đổi: 3 tấn 5 tạ = \dots tạ.
\shortans{35}
\loigiai{
35 tạ.
}
\end{ex}

\begin{ex}
Đổi: $4\text{ m}^2\ 25\text{ dm}^2 = \dots\text{ dm}^2$.
\shortans{425}
\loigiai{
425.
}
\end{ex}

\begin{ex}
Đổi: 1 phút 40 giây = \dots giây.
\shortans{100}
\loigiai{
100 giây.
}
\end{ex}

\begin{ex}
Năm 2026 thuộc thế kỉ thứ mấy? (Ghi số tự nhiên).
\shortans{21}
\loigiai{
Thế kỉ XXI (21).
}
\end{ex}

\begin{ex}
Hai đường thẳng vuông góc tạo thành góc bao nhiêu độ?
\shortans{90}
\loigiai{
$90^\circ$.
}
\end{ex}

\begin{ex}
Góc bẹt bằng bao nhiêu độ?
\shortans{180}
\loigiai{
$180^\circ$.
}
\end{ex}

\begin{ex}
Hình thoi có hai đường chéo dài 12 cm và 10 cm. Diện tích là bao nhiêu xăng-ti-mét vuông?
\shortans{60}
\loigiai{
$12 \times 10 : 2 = 60\text{ cm}^2$.
}
\end{ex}

\begin{ex}
Hình bình hành có đáy 15 cm, chiều cao 6 cm. Diện tích là bao nhiêu xăng-ti-mét vuông?
\shortans{90}
\loigiai{
$15 \times 6 = 90\text{ cm}^2$.
}
\end{ex}

\begin{ex}
Một hình chữ nhật có nửa chu vi là 30 m, chiều dài hơn chiều rộng 6 m. Chiều dài là bao nhiêu mét?
\shortans{18}
\loigiai{
$(30 + 6) : 2 = 18$ m.
}
\end{ex}

\begin{ex}
Chiều rộng hình chữ nhật trong bài toán trên là bao nhiêu mét?
\shortans{12}
\loigiai{
$(30 - 6) : 2 = 12$ m.
}
\end{ex}

\begin{ex}
Diện tích hình chữ nhật đó là bao nhiêu mét vuông?
\shortans{216}
\loigiai{
$18 \times 12 = 216\text{ m}^2$.
}
\end{ex}
